\subsection{Case of algebras without a unit element}

DEFINITION 6. --- Let A be a commutative algebra over K. A character of A is an algebra homomorphism from A to K.

The set of characters of A will be denoted by X'(A).

The zero map is an algebra homomorphism. If A has a unit element, \(e\) , a nonzero algebra morphism from A into K is unital, that is, it is a unital character of A in the sense of Definition 4: indeed, for \(\chi \in \mathsf{X}'(\mathsf{A})\), \(\chi\) is nonzero if and only if \(\chi(e) = 1\) .

We set \(\mathsf{X}(\mathsf{A})=\mathsf{X}^{\prime}(\mathsf{A})-\{0\}\) ; by the preceding, the notation is compatible with that introduced when A is unital.

If \(h\colon A\to B\) is a morphism of commutative algebras, the map \(\chi\mapsto\chi\circ h\) is a map \(\mathsf{X}^{\prime}(h)\colon\mathsf{X}^{\prime}(\mathrm{B})\to\mathsf{X}^{\prime}(\mathrm{A})\) . It maps 0 to 0. If \(k\colon B\to C\) is a morphism of commutative algebras, then one has \(\mathsf{X}^{\prime}(k\circ h)=\mathsf{X}^{\prime}(h)\circ\mathsf{X}^{\prime}(k)\) . If \(h\) is surjective, \(\mathsf{X}^{\prime}(h)\) is a bijection from \(\mathsf{X}^{\prime}(\mathrm{B})\) onto the set of characters of A vanishing on the kernel of \(h\) . Let \(\mathrm{A}_{1},\ldots,\mathrm{A}_{n}\) be commutative algebras, \(\mathrm{A}=\mathrm{A}_{1}\times\cdots\times\mathrm{A}_{n}\) and \(\pi_{i}\colon\mathrm{A}\to\mathrm{A}_{i}\) the canonical morphism; then \(\mathsf{X}^{\prime}(\pi_{i})\) is a bijection from \(\mathsf{X}^{\prime}(\mathrm{A}_{i})\) onto a subset \(\mathsf{X}_{i}^{\prime}\) of \(\mathsf{X}^{\prime}(\mathrm{A})\) , namely the set of characters of A vanishing on \(\prod_{j\neq i}\mathrm{A}_{j}\) ; we see as in no. 6 that \(\mathsf{X}^{\prime}(\mathrm{A})\) is the union

of \(X_{i}^{\prime}\) ; on the other hand, \(X_{i}^{\prime} \cap X_{j}^{\prime} = \{0\}\) for \(i \neq j\) ; in particular the \(X_{i}^{\prime} - \{0\}\) form a partition of \(X^{\prime}(A) - \{0\} = X(A)\) .

For every \(x \in A\) , let \(\mathcal{G}_{\mathrm{A}}^{\prime}(x)\) , or simply \(\mathcal{G}^{\prime}(x)\) , be the map \(\chi \mapsto \chi(x)\) from \(\mathsf{X}^{\prime}(\mathrm{A})\) to K. The map \(\mathcal{G}^{\prime}\) is a morphism from A into the algebra \(A_{1}\) of the maps \(\mathsf{X}^{\prime}(\mathrm{A}) \to \mathrm{K}\) vanishing at 0. Let B be a commutative algebra, let \(B_{1}\) be the algebra of maps \(\mathsf{X}^{\prime}(\mathrm{B}) \to \mathrm{K}\) vanishing at 0, and let h be a morphism from A to B; then \(\mathsf{X}^{\prime}(h)\) defines a morphism \(h_{1}: A_{1} \to B_{1}\) , and we have \(h_{1} \circ \mathcal{G}_{A}^{\prime} = \mathcal{G}_{B}^{\prime} \circ h\) . One denotes by \(\mathcal{G}_{\mathrm{A}}(x)\) , or simply \(\mathcal{G}(x)\) the restriction of \(\mathcal{G}_{\mathrm{A}}^{\prime}(x)\) onto \(\mathrm{X}(\mathrm{A})\) and it is called the Gelfand transform of \(x\) .

Let \(\widetilde{\mathsf{A}}\) be the unital algebra deduced from A by adjoining an identity element. By restriction, every character of \(\widetilde{\mathsf{A}}\) defines a character of A; conversely, every character of A extends uniquely to a character of \(\widetilde{\mathsf{A}}\) . This defines a canonical bijection from \(\mathsf{X}'(\mathsf{A})\) onto \(\mathsf{X}(\widetilde{\mathsf{A}})\) , by which these two sets are identified. The character 0 of A is identified with the unique character of \(\widetilde{\mathsf{A}}\) with kernel A.

If \(x\in \mathrm{A}\) and \(\chi \in \mathsf{X}'(\mathrm{A})\) , one has \(\chi (x)\in \mathrm{Sp}_{\widetilde{\mathrm{A}}}^{\sim}(x)\) , hence \(\chi (x)\in \mathrm{Sp}_{\mathrm{A}}^{\prime}(x)\) .

Lemma 3. --- The map \(\chi \mapsto \mathrm{Ker}(\chi)\) is a bijection from \(\mathsf{X}(\mathsf{A})\) onto the set of regular ideals of codimension 1 of A.

Recall (A, VIII, p.~426, def. 1) that an ideal I of A is said to be regular if the quotient algebra A/I admits a unit element.

Let us prove the lemma. On the one hand, \(\mathsf{X}(\mathsf{A})\) is identified with the set of characters of \(\widetilde{A}\) nonzero on A. On the other hand, by A, VIII, p.~428, prop. 4, the map \(I \mapsto A \cap I\) is a bijection from the set of maximal ideals of \(\widetilde{A}\) distinct from A onto the set of regular maximal ideals of A. The lemma then follows from the results of no. 6.

Suppose now that K is a topological field. We then equip \(\mathsf{X}^{\prime}(\mathsf{A})\) with the topology of pointwise convergence on A; the notation \(\mathsf{X}^{\prime}(\mathsf{A})\) will henceforth denote the topological space thus obtained. When K = R or C, we shall also call it the weak topology. For every \(x \in A\) , the function \(\mathcal{G}_{\mathrm{A}}^{\prime}(x)\) on \(\mathsf{X}^{\prime}(\mathsf{A})\) is continuous.

If h is a morphism from A to B, the map \(\mathsf{X}'(h)\colon\mathsf{X}'(\mathsf{B})\to\mathsf{X}'(\mathsf{A})\) is continuous. If h is surjective, \(\mathsf{X}'(h)\) is a homeomorphism from \(\mathsf{X}'(\mathsf{B})\) onto its image and this image is closed in \(\mathsf{X}'(\mathsf{A})\) .

Let \(A = A_{1} \times \cdots \times A_{n}\) ; with the same notations as above, \(\mathsf{X}'(\pi_{i})\) is a homeomorphism from \(\mathsf{X}'(\mathrm{A}_{i})\) onto \(X_{i}'\) and \(X_{i}'\) is closed in \(\mathsf{X}'(\mathrm{A})\) . Therefore \(X_{i}' - \{0\}\) is open in \(\mathsf{X}'(\mathrm{A})\) ; the \(\mathsf{X}'(\pi_{i})\) define a continuous map from the sum space S of the \(\mathsf{X}'(\mathrm{A}_{i})\) onto \(\mathsf{X}'(\mathrm{A})\) , and one verifies immediately that the image of a union of neighborhoods of the points \(0 \in \mathsf{X}'(\mathrm{A}_{1}), \ldots, 0 \in \mathsf{X}'(\mathrm{A}_{n})\) is a neighborhood of \(0 \in \mathsf{X}'(\mathrm{A})\) ; from all this it follows that \(\mathsf{X}'(\mathrm{A})\) is canonically identified with a quotient space of S. In particular, the space \(\mathsf{X}(\mathrm{A})\) is identified with the sum space of the \(\mathsf{X}(\mathrm{A}_{i})\) .

The canonical bijection from \(\mathsf{X}^{\prime}(\mathsf{A})\) onto \(\mathsf{X}(\tilde{\mathsf{A}})\) is a homeomorphism. Let B be a unital algebra over K and \(B^{\prime}\) the underlying algebra; then the space \(\mathsf{X}(\mathsf{B})\) is identified with the subspace \(\mathsf{X}(\mathsf{B}^{\prime})\) of \(\mathsf{X}^{\prime}(\mathsf{B}^{\prime})\) .

